\section{Mode Collapse and Verified Replay}
\label{app:theory}

For a fixed prompt, the binary-reward score estimators defined below have
mean updates along the correctness gradient, while their sampling
covariances can differ. With independently optimized categorical logits,
these updates amplify differences among equally rewarded modes, and sampling
can break an initially exact symmetry. Fixed-bank replay supplies
probability floors, restoration thresholds, and recovery bounds under the
stated update assumptions. Neural retention requires additional conditions
on complete-response likelihoods and local updates. These results do not
certify the implemented AdamW/PPO trajectories.
Table~\ref{tab:notation} defines the notation.

\subsection{Categorical abstraction and mean updates}
\label{app:theory-grpo-mean}

The categorical model assigns one independent logit to each outcome
category. Its Euclidean mean flow and the subsequent finite-step, sampled,
and neural updates have different assumptions; a neural solution mode need
not have an independently optimized logit.

\begin{assumption}[Categorical mean-flow model]
\label{ass:categorical-mean-flow}
Fix one isolated prompt and a finite category set
$\mathcal A=\mathcal C\mathbin{\dot\cup}\mathcal N$ with $m\ge2$ correct
solution modes and $n\ge1$ incorrect categories, and set
$p_a=\operatorname{softmax}(z)_a$, $P=\sum_{c\in\mathcal C}p_c$, and
$q_c=p_c/P$, so that $P$ measures correctness while $q$ describes how that
correct probability is allocated among the individual correct modes.
\begin{enumerate}
\item[(A1)] \emph{Independent logits.} The policy carries one independently
  optimized logit for every category in $\mathcal A$, including one separate
  logit for each canonical correct mode.
\item[(A2)] \emph{Common length normalization.} Every response uses the same
  fixed length-normalization factor, absorbed into the positive coefficient
  of the group estimator.
\item[(A3)] \emph{Euclidean mean flow.} Updates are the infinitesimal expected
  updates in these logits, taken without optimizer preconditioning and without
  momentum carried between steps.
\item[(A4)] \emph{No competing regularizer.} Explicit reference KL is zero and
  no entropy term is present.
\item[(A5)] \emph{Nondegenerate start.} Logits are finite initially,
  $0<P(0)<1$, and the group size is $G\ge2$.
\end{enumerate}
\end{assumption}

Assumption~\ref{ass:categorical-mean-flow} includes all five conditions.
AdamW's preconditioning and momentum violate (A3); finite, repeated PPO
updates are outside its infinitesimal on-policy limit. Response-specific
length normalization violates (A2), while reference KL and entropy
regularization violate (A4).

A language model induces probabilities over canonical outcomes, but those
probabilities alone do not determine its update geometry. Aggregating several
response strings into a mode does not generally commute with gradient flow
in response logits. The independent-category model therefore requires an
additional assumption beyond a neural policy and its validator equivalence
classes.

Draw a group $Z_1,\ldots,Z_G\stackrel{\mathrm{iid}}{\sim}p$ with $G\ge2$,
write $R_i=\1[Z_i\in\mathcal C]$, and let $R=\sum_iR_i$. At the on-policy
point where the PPO ratio equals one, consider the common-length estimator
\begin{equation}
  \widehat g_G
  =\frac1G\sum_{i=1}^G
   w(R)\left(R_i-\frac{R}{G}\right)\nabla_z\log p_{Z_i}.
  \label{eq:group-estimator}
\end{equation}
For the Dr.GRPO abstraction \citep{liu2025understanding}, $w(R)=1$ after
absorbing the fixed positive length-normalization factor. Reward-standardized
GRPO \citep{shao2024deepseekmath} with fixed common response-length
normalization takes $w(R)$ to be the reciprocal group reward standard
deviation on mixed groups (with a fixed numerical stabilizer if used).
Standard variable-length GRPO also weights each response by its own inverse
length; that response-specific factor cannot generally be absorbed into
$w(R)$ and is outside this equivalence.
In either included case, $0<w(\ell)<\infty$ for
$\ell\in\{1,\ldots,G-1\}$. The entire update is zero on all-equal
groups, without evaluating an undefined reciprocal standard deviation;
likewise, $w(R)R(G-R)$ below is defined as zero when $R=0$ or $R=G$. All
advantages and normalization weights are detached in the score estimator.

\begin{lemma}[A group baseline does not remove frequency amplification]
\enforcehalffinalproseline
\label{lem:grpo-mean}
For $0<P<1$, the expected update in Equation~\eqref{eq:group-estimator} is
\begin{equation}
  \E[\widehat g_G]
  =c_G(P)\nabla_z P,
  \qquad
  c_G(P)=
  \frac{\E\!\big[w(R)R(G-R)\big]}{G^2P(1-P)}>0.
  \label{eq:expected-group-gradient}
\end{equation}
For the Dr.GRPO abstraction, $c_G(P)=(G-1)/G$. Group centering and reward
standardization rescale the expected gradient without changing its direction
under these assumptions.
\end{lemma}

\begin{proof}
\enforcehalffinalproseline
Condition on the complete binary reward vector, whose sum is $R$. The score
identities give
\[
  \E[\nabla_z\log p_{Z_i}\mid R_i=1]=\nabla_z\log P,
  \qquad
  \E[\nabla_z\log p_{Z_i}\mid R_i=0]
  =\nabla_z\log(1-P).
\]
There are $R$ correct rows with centered reward $(G-R)/G$ and $G-R$
incorrect rows with centered reward $-R/G$. Therefore the conditional
expectation of the unaveraged sum in~\eqref{eq:group-estimator} is
\[
  w(R)\frac{R(G-R)}{G}
  \left(\nabla_z\log P-\nabla_z\log(1-P)\right)
  =w(R)\frac{R(G-R)}{GP(1-P)}\nabla_zP.
\]
The outer factor $1/G$ and expectation over
$R\sim\operatorname{Binomial}(G,P)$ yield~\eqref{eq:expected-group-gradient}.
For $w\equiv1$,
$\E[R(G-R)]=G(G-1)P(1-P)$, proving the Dr.GRPO value.
\end{proof}

The binomial identity also gives the denominator-free expression
\[
 c_G(P)=\frac{G-1}{G}\sum_{j=0}^{G-2}\binom{G-2}{j}
 w(j+1)P^j(1-P)^{G-2-j}.
\]
The binomial weights sum to one, so $c_G$ lies between
$(G-1)\min_j w(j+1)/G$ and $(G-1)\max_j w(j+1)/G$.
Thus $c_G$ extends smoothly and positively to $P=0,1$.

The practical MaxRL estimator has the same mean-gradient direction
\citep[Section~4.3 and App.~D, Theorem~5]{tajwar2026maxrl}. Its all-failure
convention determines the finite-group coefficient.

\begin{lemma}[Binary MaxRL has the same correctness-gradient direction]
\enforcehalffinalproseline
\label{lem:maxrl-mean}
Under Assumption~\ref{ass:categorical-mean-flow}, let MaxRL assign
$A_i^{\mathrm{MaxRL}}=GR_i/R-1$ when $R>0$ and zero when $R=0$, as in
Section~\ref{sec:method}. For
$\widehat g_G^{\mathrm{MaxRL}}=G^{-1}\sum_i
 A_i^{\mathrm{MaxRL}}\nabla_z\log p_{Z_i}$ at the on-policy point, the
expected update is
\begin{equation}
  \E[\widehat g_G^{\mathrm{MaxRL}}]
  =c_G^{\mathrm{MaxRL}}(P)\nabla_zP,
  \qquad
  c_G^{\mathrm{MaxRL}}(P)
  =\frac{1-(1-P)^{G-1}}{P}>0,
  \label{eq:maxrl-expected-gradient}
\end{equation}
with continuous endpoint values
$c_G^{\mathrm{MaxRL}}(0)=G-1$ and
$c_G^{\mathrm{MaxRL}}(1)=1$.
\end{lemma}

\begin{proof}
\enforcehalffinalproseline
With $N=G$ and pass rate $P$, the estimator in
\citet[Theorem~5]{tajwar2026maxrl} equals
$\widehat g_G^{\mathrm{MaxRL}}$. Its expected coefficient is
$\sum_{j=0}^{G-2}(1-P)^j$, whose geometric sum gives
Equation~\eqref{eq:maxrl-expected-gradient} and its endpoint values.
The identity requires on-policy independent draws and the score identity.
\end{proof}

Integrating this coefficient gives the potential
\begin{equation}
 \Psi_{\mathrm{MaxRL}}(P)
 :=\int_0^P c_G^{\mathrm{MaxRL}}(s)\,ds
 =\sum_{k=1}^{G-1}\frac{1-(1-P)^k}{k}.
 \label{eq:maxrl-potential-finite-sum}
\end{equation}
In particular, $\Psi_{\mathrm{MaxRL}}(1)=\sum_{k=1}^{G-1}1/k$.
For binary rewards, best@\(k\) equals pass@\(k\), so this is a harmonic
mixture of sampling objectives for $k=1,\ldots,G-1$. The upper limit is
$G-1$ because the implemented centered estimator sets the entire update
to zero on all-failure groups. Subtracting an unconditional score baseline
even on those groups instead gives the order-$G$ estimator. These are exact
finite-group expectations at the on-policy point; subsequent clipping and
adaptive optimization can alter the expected update.


\subsection{Winner-take-all collapse in the categorical mean flow}
\label{app:theory-collapse}

Consider $\dot z=c_G(P)\nabla_zP$, with any of the smooth positive
coefficients established above: Dr.GRPO, common-length reward-standardized
GRPO, or the specified centered binary MaxRL estimator. This infinitesimal
expected-update model has
no sampling noise and uses independently optimized Euclidean logits. The
conditional dynamics follow the frequency-dependent replicator system
with fitness $f_c(q)=q_c$
\citep{hofbauer1998evolutionary,harper2009informationgeometry}.
Related analyses of collapse among equally rewarded outcomes include
\citet{sinha2026expected,lochab2026ucpo}.

\begin{theorem}[Winner-take-all collapse in the categorical mean flow]
\enforcehalffinalproseline
\label{thm:grpo-collapse}
Under Assumption~\ref{ass:categorical-mean-flow}, if the conditional correct-mode
distribution $q(0)$ has a unique largest coordinate $c^\star$, then
\[
  P(t)\longrightarrow1,
  \qquad
  q_{c^\star}(t)\longrightarrow1,
  \qquad
  q_c(t)\longrightarrow0\quad(c\ne c^\star).
\]
Thus these categorical mean flows converge to a single correct solution
mode for generic initial conditions, even though every mode in $\mathcal C$
has exactly the same reward.
\end{theorem}

\begin{remark}
Theorem~\ref{thm:grpo-collapse} is asymptotic and requires (A1)--(A5).
Finite-time support loss and convergence of the stochastic neural optimizers
used in the experiments are outside its scope.
\end{remark}

\begin{proof}
\enforcehalffinalproseline
The coefficient $c_G$ is smooth and bounded on $[0,1]$, and the softmax
reward gradient is bounded. Thus the logit flow exists for all finite times
and cannot develop infinite logits at finite time.
For a correct mode $c$ and incorrect category $a$,
\[
  \dot z_c=c_G(P)p_c(1-P),
  \qquad
  \dot z_a=-c_G(P)p_aP.
\]
Also $\dot P=c_G(P)\lVert\nabla_zP\rVert_2^2>0$. If $P$ converged to an
interior value, continuity and positivity of $c_G$, together with
\[
 \lVert\nabla_zP\rVert_2^2
 \ge P^2(1-P)^2\left(\frac1m+\frac1n\right),
\]
would keep $\dot P$ bounded away from zero, a contradiction. Hence $P\to1$.

Define the effective optimization time
\[
  \tau(t)=\int_0^t c_G(P(s))P(s)(1-P(s))\,ds.
\]
If $\tau(\infty)<\infty$, every correct and incorrect logit would have finite
total variation, because the absolute velocity of each is at most $\dot\tau$.
All logits would then converge to finite values, contradicting $P\to1$.
Therefore $\tau(t)\to\infty$, leaving infinite effective time for the
correct-simplex dynamics below.

On the correct simplex, changing from $t$ to $\tau$ gives the replicator flow
\begin{equation}
  \frac{dq_c}{d\tau}
  =q_c\left(q_c-\sum_{d\in\mathcal C}q_d^2\right),
  \qquad
  \frac{d}{d\tau}\log\frac{q_c}{q_d}=q_c-q_d.
  \label{eq:grpo-replicator}
\end{equation}
The initially unique maximum remains the unique maximum. For $c\ne c^\star$,
set $\xi_c=q_c/q_{c^\star}$, so $\xi_c(0)<1$ and $\xi_c$ is nonincreasing.
Since $q_{c^\star}\ge1/m$,
\[
 \frac{d\log\xi_c}{d\tau}
 =-q_{c^\star}(1-\xi_c)
 \le-\frac{1-\xi_c(0)}{m}.
\]
Hence $\xi_c(\tau)\le\xi_c(0)\exp[-(1-\xi_c(0))\tau/m]\to0$.
Every competing mode vanishes relative to $c^\star$; normalization then
yields $q_{c^\star}\to1$, establishing the claimed single-mode limit.
\end{proof}

If several modes tie for the largest initial probability, symmetry preserves
the tie; the same ratio argument eliminates strictly smaller coordinates and
leaves the tied maxima uniform. An exactly uniform correct distribution is
therefore a special initial condition. The theorem concerns limiting
concentration: finite softmax logits already give positive probabilities at
every finite time, so finite-time positivity alone is not a retention result.

\subsection{Conditional trajectories of outcome-binary mean flows}
\label{app:theory-objective-inertness}

Under Assumption~\ref{ass:categorical-mean-flow}, every outcome-binary
estimator with a positive correctness coefficient induces the same
conditional trajectory in $q$. Its advantage changes only the rate at which
that trajectory is traversed. Conditioning on correctness averages over all
correct modes, so a reward-dependent advantage cannot distinguish their
identities.

\paragraph{Proof of Theorem~\ref{thm:objective-inertness}.}

\begin{proof}
\enforcehalffinalproseline
The score identities used in Lemma~\ref{lem:grpo-mean} give
$\E[\nabla_z\log p_{Z_i}\mid R_i=1]=\nabla_z\log P$ and
$\E[\nabla_z\log p_{Z_i}\mid R_i=0]=\nabla_z\log(1-P)$.
Independence of the draws means that conditioning also on the other rows'
rewards does not change either conditional mean. Since $a(R_i,R)$ depends
only on the binary rewards, conditioning on their complete vector and
summing the $R$ successful and $G-R$ failed rows gives
\[
  \E[\widehat g_a]
  =\frac1G\Big(\E[R\,a(1,R)]\nabla_z\log P
   +\E[(G-R)\,a(0,R)]\nabla_z\log(1-P)\Big).
\]
Substituting $\nabla_z\log P=\nabla_zP/P$ and
$\nabla_z\log(1-P)=-\nabla_zP/(1-P)$ yields
Equation~\eqref{eq:objective-inertness}. Hence
$\frac{d}{dt}\log(q_c/q_d)=\dot z_c-\dot z_d=c_a(P)P(1-P)(q_c-q_d)$.
For $c_a>0$, setting $d\tau=c_a(P)P(1-P)\,dt$ gives
Equation~\eqref{eq:grpo-replicator}.
\end{proof}

For a fixed reward-dependent advantage, the binomial identities
$R\binom GR=G\binom{G-1}{R-1}$ and
$(G-R)\binom GR=G\binom{G-1}{R}$ give
\begin{equation}
 c_a(P)=\sum_{j=0}^{G-1}\binom{G-1}{j}P^j(1-P)^{G-1-j}
 \big[a(1,j+1)-a(0,j)\big],
 \label{eq:inertness-polynomial}
\end{equation}
a polynomial in $P$ of degree at most $G-1$. Finite advantage values
therefore give a smooth, bounded coefficient on $[0,1]$.

The coefficients in this Bernstein expansion are
\begin{equation}
 d_j:=a(1,j+1)-a(0,j),\qquad j=0,\ldots,G-1 .
 \label{eq:leave-one-out-gap}
\end{equation}
With the other $G-1$ rows fixed at $j$ successes, $d_j$ is the change in a
row's advantage when its reward changes from failure to success. The
Bernstein basis is nonnegative on $[0,1]$ and sums to one, so $c_a(P)$ lies
between $\min_jd_j$ and $\max_jd_j$. Each basis function is strictly positive
on $(0,1)$; hence nonnegative gaps, with at least one positive, imply
$c_a(P)>0$ throughout the interior. This condition is sufficient:
mixed-sign gaps alone do not determine the sign of $c_a$.

For Dr.GRPO, every $d_j=(G-1)/G$, giving the constant coefficient in
Lemma~\ref{lem:grpo-mean}. For the centered MaxRL estimator,
$d_0=G-1$ and $d_j=G/(j+1)$ for $1\le j\le G-1$.
In particular, $a(1,G)=0$ and $a(0,G-1)=-1$ give $d_{G-1}=1$.
Substitution into Equation~\eqref{eq:inertness-polynomial} gives
$[1-(1-P)^{G-1}]/P$, as in Lemma~\ref{lem:maxrl-mean}.

\begin{corollary}[Collapse for correctness-improving outcome-binary mean flows]
\enforcehalffinalproseline
\label{cor:class-collapse}
Under Assumption~\ref{ass:categorical-mean-flow}, let $a(R_i,R)$ be a fixed,
bounded advantage depending only on the binary rewards. Suppose the gaps
in Equation~\eqref{eq:leave-one-out-gap} satisfy $d_j\ge0$ for every $j$,
with at least one strictly positive gap. If $q(0)$ has a unique largest
coordinate $c^\star$, then
\[
 P(t)\longrightarrow1,\qquad
 q_{c^\star}(t)\longrightarrow1,\qquad
 q_c(t)\longrightarrow0\quad(c\ne c^\star).
\]
The surviving correct mode is the one with the largest initial probability;
the choice of $a$ changes the rate along the conditional trajectory.
\end{corollary}

\begin{proof}
\enforcehalffinalproseline
By Theorem~\ref{thm:objective-inertness}, the flow is
$\dot z=c_a(P)\nabla_zP$. Equation~\eqref{eq:inertness-polynomial} makes
$c_a$ smooth and bounded on $[0,1]$. If $d_{j_0}>0$, the nonnegative-gap
hypothesis gives
\[
 c_a(P)\ge\binom{G-1}{j_0}P^{j_0}(1-P)^{G-1-j_0}d_{j_0}>0
 \qquad(0<P<1).
\]
These are the coefficient properties used in the proof of
Theorem~\ref{thm:grpo-collapse}: the logit flow exists for all finite times,
$P\to1$, and the effective time $\tau\to\infty$. The conditional
replicator equation~\eqref{eq:grpo-replicator} then eliminates all coordinates
strictly below the initial maximum. It contains neither the advantage
coefficients nor the rollout group size, so the surviving mode depends only
on $q(0)$.
\end{proof}

The nonnegative-gap condition is sufficient rather than necessary. An
estimator with mixed-sign gaps also has the same limiting behavior if
$c_a(P)>0$ throughout $(0,1)$. An interior zero of $c_a$ is a stationary
correctness level for the mean flow, while a negative coefficient gives
$\dot P=c_a(P)\lVert\nabla_zP\rVert_2^2<0$ at an interior policy.
The collapse conclusion requires positivity and an initial imbalance under
Assumption~\ref{ass:categorical-mean-flow}; it does not describe shared-prompt
or sampled neural updates.

\begin{remark}[Group size in the deterministic conditional flow]
\label{rem:group-size-inert}
For each $G\ge2$, group size enters the mean flow through $c_a(P)$.
When this coefficient is positive, the time change in
Theorem~\ref{thm:objective-inertness} gives the same conditional replicator
equation. Increasing $G$ can change the speed of the mean flow but cannot
change the surviving mode under the conditions of
Corollary~\ref{cor:class-collapse}.
Group size also changes the probability of a mixed-reward group in
Lemma~\ref{lem:replay-gradient-availability}, which is small near both
$P=0$ and $P=1$, and the MaxRL coefficient in
Corollary~\ref{cor:maxrl-starvation}. These properties concern total
correctness $P$; they do not provide a probability floor for an individual
correct mode. Gradient covariance and finite-sample symmetry breaking can depend on $G$,
so sampled updates need not share this equivalence.
\end{remark}

\begin{corollary}[MaxRL amplification of the correctness gradient]
\enforcehalffinalproseline
\label{cor:maxrl-starvation}
For the coefficients of Lemmas~\ref{lem:grpo-mean} and~\ref{lem:maxrl-mean},
\[
 \frac{c_G^{\mathrm{MaxRL}}(P)}{c_G^{\mathrm{Dr.GRPO}}}
 =\frac{G\big[1-(1-P)^{G-1}\big]}{(G-1)P}.
\]
This ratio is nonincreasing in $P$, with limits $G$ as $P\to0$ and
$G/(G-1)$ as $P\to1$. It is strictly decreasing for $G>2$ and constant
at $2$ for $G=2$. At the same policy, MaxRL therefore multiplies the
Dr.GRPO correctness-gradient coefficient by a factor approaching $G$ when
correctness is small. The probability of a nonzero fresh task component
remains bounded by the mixed-group probability in
Lemma~\ref{lem:replay-gradient-availability} for both estimators.
\end{corollary}

\begin{proof}
\enforcehalffinalproseline
Divide $c_G^{\mathrm{MaxRL}}(P)=[1-(1-P)^{G-1}]/P$ by $(G-1)/G$.
The finite-series identity
$c_G^{\mathrm{MaxRL}}(P)=\sum_{j=0}^{G-2}(1-P)^j$
gives the endpoint limits and monotonicity. For $G=2$ the sum is the
constant $1$; for $G>2$ it contains the strictly decreasing term $1-P$.
\end{proof}

\subsection{Neural score identities and the limits of conditional inertness}
\label{app:neural-binary-inertness}

Expected binary reward equals the total correct mass; conditional diversity
depends on its allocation among solution modes.
\citet[Theorem~3.1, v1]{sinha2026expected} derive the expected-return
log-ratio dynamics for softmax outcome logits, and
\citet[Section~4.1, v1]{lochab2026ucpo} identify binary reward's
indifference to allocation within the correct set. For differentiable neural
policies, outcome-binary group estimators have per-prompt means proportional
to the correctness gradient but can differ in covariance. Finite steps,
shared prompts, and response-dependent weighting limit the resulting
trajectory equivalence.

\begin{theorem}[Per-prompt neural mean and covariance]
\enforcehalffinalproseline
\label{thm:neural-binary-moments}
Fix a prompt and a differentiable policy $\pi_\theta$ on a finite or countable
response space with common positive support, where $\theta$ lies in an open
subset of $\mathbb R^d$. Let $\mathcal C$ be a fixed verifier-positive event,
$P=\Pr_\theta(Y\in\mathcal C)\in(0,1)$, and
$s(Y)=\nabla_\theta\log\pi_\theta(Y)$. Assume differentiation may be
interchanged with the probability sums over the full support and over
$\mathcal C$, and that $\E\|s(Y)\|_2^2<\infty$. For $G$ independent
responses, write $R_i=\1[Y_i\in\mathcal C]$, $R=\sum_iR_i$, and let
$a:\{0,1\}\times\{0,\ldots,G\}\to\mathbb R$ be any finite, detached
advantage. Define
\[
 \widehat g_a=\frac1G\sum_i a(R_i,R)s(Y_i),\qquad
 v=\nabla_\theta P,\qquad
 \Sigma_r=\operatorname{Cov}\!\left(s(Y)\mid\1[Y\in\mathcal C]=r\right).
\]
With
\[
 \beta_r=\frac1G\left(\frac{r\,a(1,r)}P
          -\frac{(G-r)a(0,r)}{1-P}\right),
 \qquad c_a(P)=\E[\beta_R],
\]
where $R\sim\operatorname{Binomial}(G,P)$, the exact identities are
\begin{align}
 \E[\widehat g_a]&=c_a(P)v,
 \label{eq:neural-binary-mean}\\
 \operatorname{Cov}(\widehat g_a)
 &=\operatorname{Var}(\beta_R)vv^\top
   +\frac1{G^2}\E\!\left[
       R\,a(1,R)^2\Sigma_1+(G-R)a(0,R)^2\Sigma_0\right].
 \label{eq:neural-binary-covariance}
\end{align}
In particular, $c_a$ is the coefficient in
Equation~\eqref{eq:objective-inertness}, without an independent-logit
assumption.
\end{theorem}

\begin{proof}
\enforcehalffinalproseline
The conditional-score identity
\citep[Theorem~1]{tajwar2026maxrl} (also
\citealp[Lemma~B.1]{avspo2026}) applied to the correct and incorrect events gives
\[
 \E[s(Y)\mid Y\in\mathcal C]=v/P,
 \qquad
 \E[s(Y)\mid Y\notin\mathcal C]=-v/(1-P).
\]
Conditioned on the complete reward vector, the response scores remain
independent. The conditional mean of $\widehat g_a$ is $\beta_Rv$ and
its conditional covariance is
$G^{-2}[R a(1,R)^2\Sigma_1+(G-R)a(0,R)^2\Sigma_0]$.
Taking expectations proves Equation~\eqref{eq:neural-binary-mean}; the law
of total covariance proves Equation~\eqref{eq:neural-binary-covariance}.
\end{proof}

For one prompt, a common locally Lipschitz preconditioner $M(\theta)$ and
positive coefficients $c_a(P)$ give deterministic mean fields
$\dot\theta=c_a(P)M(\theta)\nabla_\theta P$ that trace the same orbit
under a positive change of time, wherever the flows exist uniquely.
This is a statement about the mean direction. Equation~\eqref{eq:neural-binary-covariance}
shows that outcome-binary estimators can have different sampling covariance,
including after their mean magnitudes are matched.

\begin{corollary}[A finite-step comparison in neural parameters]
\enforcehalffinalproseline
\label{cor:neural-binary-finite-step}
Under Theorem~\ref{thm:neural-binary-moments}, suppose $c_a(P)>0$ and take
$\theta^+=\theta+\eta\widehat g_a/c_a(P)$ with $\eta>0$. Let $f$ be a
twice differentiable scalar observable with
$\|\nabla^2f\|_{\mathrm{op}}\le L_f$ along every update segment, which
is assumed to remain in its domain. Then, conditional on the current policy,
\begin{equation}
 \left|\E[f(\theta^+)-f(\theta)]
       -\eta\nabla f(\theta)^\top v\right|
 \le\frac{L_f\eta^2}{2}
 \left(\|v\|_2^2+
       \frac{\operatorname{tr}\operatorname{Cov}(\widehat g_a)}{c_a(P)^2}
 \right).
 \label{eq:neural-binary-finite-step}
\end{equation}
\end{corollary}

\begin{proof}
\enforcehalffinalproseline
Taylor's theorem bounds the remainder after the linear term by
$L_f\eta^2\|\widehat g_a\|_2^2/(2c_a(P)^2)$. Take expectations and use
Equation~\eqref{eq:neural-binary-mean} together with
$\E\|\widehat g_a\|_2^2
=c_a(P)^2\|v\|_2^2+\operatorname{tr}\operatorname{Cov}(\widehat g_a)$.
\end{proof}

For two estimators satisfying the corollary, division by their respective
coefficients $c_a(P)$ gives the same mean step $\eta v$; the training updates
do not use this normalization. Their expected changes in $f$ then differ by
at most the sum of the two second-order bounds. Conditional solution mode
probabilities or \pmd{} may be used as $f$ where the stated smoothness bounds
hold. Mean-direction equivalence is
therefore a first-order statement for stochastic finite updates; it does
not determine their trajectories or eventual surviving solution modes.
The bound assumes fresh on-policy score updates and does not cover Adam's
moment estimates or subsequent PPO epochs.

\begin{remark}[Shared prompts need not follow the same trajectory]
\label{rem:neural-binary-shared-prompts}
For a fixed prompt-sampling distribution $\mu_x$ and shared parameters,
the aggregate mean, assuming integrable updates, is
\begin{equation}
 \E[\widehat g_a]
 =\sum_x\mu_x c_a(P_x)\nabla_\theta P_x.
 \label{eq:neural-binary-prompt-sum}
\end{equation}
Different objectives generally assign different relative coefficients to
the prompt gradients, so the per-prompt identity does not imply a common
aggregate orbit. An explicit example also shows a difference in conditional
diversity. Let $\theta=(u,v)$, $\sigma(t)=(1+e^{-t})^{-1}$, and sample two
prompts $A,B$ equally. Set
\[
 P_A=\sigma(u+v),\qquad P_B=\sigma(u-v),\qquad
 q_A=(\sigma(v),1-\sigma(v)).
\]
For prompt $A$, the probabilities of two correct responses and one incorrect
response are $(P_A\sigma(v),P_A[1-\sigma(v)],1-P_A)$; for $B$, use
$(P_B/2,P_B/2,1-P_B)$. At $(u,v)=(0,\log3)$,
$P_A=3/4$, $P_B=1/4$, and
$\nabla P_A=(3/16,3/16)$, $\nabla P_B=(3/16,-3/16)$.
For $G=3$, Lemmas~\ref{lem:grpo-mean} and~\ref{lem:maxrl-mean} give
$c_{\mathrm{Dr}}=2/3$ and $c_{\mathrm{Max}}(P)=2-P$, hence
\[
 \dot\theta_{\mathrm{Dr}}=(1/8,0),\qquad
 \dot\theta_{\mathrm{Max}}=(9/32,-3/64).
\]
Prompt $A$ has conditional disagreement
$D_A=1-\sum_cq_{A,c}^2=2\sigma(v)[1-\sigma(v)]$, so direct
differentiation gives
\[
 \dot D_{A,\mathrm{Dr}}=0,\qquad
 \dot D_{A,\mathrm{Max}}=9/1024>0.
\]
Both prompts' correctness increases under both mean fields, since
$\dot P_A=(3/16)(\dot u+\dot v)>0$ and
$\dot P_B=(3/16)(\dot u-\dot v)>0$ at this point. Changing a binary
objective can therefore change conditional diversity through shared prompt
gradients, despite the exact per-prompt identity.
\end{remark}

\begin{proposition}[The residual from response-dependent normalization]
\enforcehalffinalproseline
\label{prop:neural-binary-length-residual}
Under Theorem~\ref{thm:neural-binary-moments}, multiply each response's score
by a detached function $\omega(Y)$ with finite conditional second moments;
inverse response length is the case $\omega(Y)=1/L(Y)$. Define
\[
 \widehat g_{a,\omega}
 =\frac1G\sum_i a(R_i,R)\omega(Y_i)s(Y_i),\qquad
 A_1=\frac{\E[R a(1,R)]}{G},\quad
 A_0=\frac{\E[(G-R)a(0,R)]}{G},
\]
and, conditioning on $\1[Y\in\mathcal C]=r$, let
$\bar\omega_r=\E[\omega(Y)\mid r]$,
$K_r=\operatorname{Cov}(\omega(Y),s(Y)\mid r)$, and
$\sigma_{\omega,r}^2=\operatorname{Var}(\omega(Y)\mid r)$.
Then
\begin{align}
 \E[\widehat g_{a,\omega}]
 &=\widetilde c_a\nabla_\theta P+A_1K_1+A_0K_0,
 &\widetilde c_a&=\frac{A_1\bar\omega_1}{P}
                  -\frac{A_0\bar\omega_0}{1-P},
 \label{eq:neural-binary-length-residual}\\
 \|A_1K_1+A_0K_0\|_2
 &\le\sum_{r=0}^1|A_r|\sigma_{\omega,r}
                 \sqrt{\operatorname{tr}\Sigma_r}.
 \label{eq:neural-binary-length-bound}
\end{align}
\end{proposition}

\begin{proof}
\enforcehalffinalproseline
Conditioning on the reward vector gives
$\E[\widehat g_{a,\omega}]
=A_1\E[\omega s\mid1]+A_0\E[\omega s\mid0]$.
For each outcome, $\E[\omega s\mid r]
=\bar\omega_r\E[s\mid r]+K_r$; substituting the conditional score
means proves Equation~\eqref{eq:neural-binary-length-residual}.
The vector Cauchy--Schwarz inequality gives
$\|K_r\|_2\le\sigma_{\omega,r}\sqrt{\operatorname{tr}\Sigma_r}$;
the triangle inequality completes the bound.
\end{proof}

The residual is a within-outcome covariance between response weighting and
the policy score. For inverse-length weighting, a common response length
makes it zero; equal mean lengths for correct and incorrect responses need
not. Response-dependent weights need not be functions of binary outcomes
alone. IPS uses outcome-frequency weights
\citep[Section~4.4, v1]{sinha2026expected}, and UCPO adds a
conditional-uniformity term \citep[Section~6, v1]{lochab2026ucpo}; these
interventions are outside the outcome-binary class analyzed above.

The verified-key diversity metric is a specialization of SetPO's
distributional functional. Its influence gives an exact metric gradient
for shared neural parameters and fixed verified execution keys.

\begin{corollary}[Verified-key specialization of SetPO]
\enforcehalffinalproseline
\label{cor:pcmd-setpo-neural}
Fix a prompt, a verifier-positive response event $\mathcal C$, and a
deterministic key map $V:\mathcal C\to\{1,\ldots,m\}$, all independent
of $\theta$. Let $\pi_\theta$ be differentiable with common positive support,
$P_\theta=\Pr_\theta(\mathcal C)>0$, and
$Q_\theta=\pi_\theta(\cdot\mid\mathcal C)$. Assume the score
$s_\theta(Y)=\nabla_\theta\log\pi_\theta(Y)$ is integrable under $Q_\theta$
and differentiation commutes with the response sums over each key.
Write $q_c=Q_\theta(V(Y)=c)$ and $S_2=\sum_cq_c^2$.
For $k(y,y')=\1[V(y)=V(y')]$ and $g(u)=1-u$, the functional
$\mathcal F(Q)=\E_{Y\sim Q}[g(\E_{Y'\sim Q}k(Y,Y'))]$ is exactly
$\pmd(q)$. The influence formula of
\citet[Theorem~4.2, v1]{li2026setpo} specializes, for $y\in\mathcal C$, to
\begin{align}
 \left.\frac{d}{d\varepsilon}
  \mathcal F((1-\varepsilon)Q_\theta+\varepsilon\delta_y)
  \right|_{\varepsilon=0}
 &=2(S_2-q_{V(y)}),\label{eq:pcmd-setpo-influence}\\
 \nabla_\theta\pmd(q)
 &=2\E_{Y\sim Q_\theta}
       [(S_2-q_{V(Y)})s_\theta(Y)].
 \label{eq:pcmd-neural-gradient}
\end{align}
\end{corollary}
\begin{proof}
\enforcehalffinalproseline
The equality kernel and $g$ satisfy SetPO's Assumption~4.1, and its local
mass is $q_{V(y)}$. Substitution into its influence formula proves the first
identity. The influence has mean zero under $Q_\theta$, so the
$-\nabla_\theta\log P_\theta$ term in the conditional score cancels,
giving the second identity.
\end{proof}
These identities permit shared neural parameters. Along an update, the
sign of the diversity change depends on its alignment with this metric
gradient; the influence formula alone supplies no monotonicity guarantee
for a given training direction.

\subsection{Diversity at finite steps and sampling budgets}
\label{app:finite-time-diversity}

Softmax outcome logits amplify probability differences among equally rewarded
solutions \citep[Theorem~3.1 and App.~A.1]{sinha2026expected}, and
sampling imbalances can reinforce this effect
\citep[Section~4.2]{lochab2026ucpo}. Under the categorical model,
\pmd{} decreases from every nonuniform correct-mode distribution along
both the mean flow and finite expected-gradient steps. Sampled updates can
also break an initially uniform distribution's symmetry. The bounds below
relate concentration to the probability of observing a competing solution
mode within a finite inference budget.

\begin{proposition}[Exact diversity drift and matched correctness]
\enforcehalffinalproseline
\label{prop:pcmd-drift}
Under Assumption~\ref{ass:categorical-mean-flow}, consider
$\dot z=c(P)\nabla_zP$ with $c(P)>0$. Write
$D(q)=1-\sum_cq_c^2$, $S_2=\sum_cq_c^2$, and
$V(q)=\sum_cq_c^3-S_2^2=\operatorname{Var}_{C\sim q}(q_C)$.
For incorrect conditional probabilities $r_a=p_a/(1-P)$, set
$T_2=\sum_ar_a^2$. In effective time $d\tau=c(P)P(1-P)dt$,
\begin{equation}
 \frac{dD}{d\tau}=-2V(q),
 \qquad
 \frac{dD}{d\log[P/(1-P)]}
 =-\frac{2V(q)}{S_2+T_2}.
 \label{eq:pcmd-drift}
\end{equation}
Thus diversity strictly decreases at every nonuniform interior $q$.
Positive-coefficient mean flows from the same initial logits follow the same
$D$-versus-$P$ curve, independently of their chosen outcome-binary advantage or
group size.
\end{proposition}

\begin{proof}
\enforcehalffinalproseline
For the replicator flow~\eqref{eq:grpo-replicator}, the potential
$U(q)=S_2/2$ has fitness $\partial U/\partial q_c=q_c$.
The fitness-variance identity of
\citet[Theorems~1--2]{harper2009informationgeometry} gives
$dU/d\tau=V(q)$, hence $dD/d\tau=-2V(q)$.
The variance vanishes exactly when the
positive coordinates of $q$ are equal. Also,
\[
 \dot P=c(P)\|\nabla_zP\|_2^2
       =c(P)P^2(1-P)^2(S_2+T_2),
\]
so $d\log[P/(1-P)]/d\tau=S_2+T_2>0$, proving the second identity.
Finally, the full logit equation in effective time is
$dz_c/d\tau=q_c$ and $dz_a/d\tau=-r_a$; it contains no $c$.
The same initial logits therefore give the same full trajectory after the
time change, and hence the same curve at matched correctness for every positive-coefficient outcome-binary estimator considered here.
\end{proof}

\begin{theorem}[Collapse under finite expected-gradient steps]
\enforcehalffinalproseline
\label{thm:finite-expected-steps}
Retain (A1), (A2), (A4), and (A5) of
Assumption~\ref{ass:categorical-mean-flow}, and replace (A3) by
\begin{equation}
 z_{k+1}=z_k+\eta_k c(P_k)\nabla_zP_k,
 \qquad 0<\eta_k<\infty,\quad \sum_{k=0}^\infty\eta_k=\infty,
 \label{eq:finite-expected-update}
\end{equation}
where $c$ is continuous and positive on $(0,1)$. Each step strictly decreases
$D(q_k)$ unless $q_k$ is uniform. If $q_0$ has a unique largest coordinate
$c^\star$, then $P_k\to1$ and $q_{c^\star,k}\to1$.
No upper bound on the finite step sizes is required.
\end{theorem}

\begin{proof}
\enforcehalffinalproseline
Let $h_k=\eta_kc(P_k)P_k(1-P_k)>0$. The exact updates are
$z_c^+=z_c+h_kq_c$ and $z_a^+=z_a-h_kr_a$, so
\begin{equation}
 q_c^+=\frac{q_c e^{h_kq_c}}{\sum_dq_d e^{h_kq_d}}.
 \label{eq:finite-correct-tilt}
\end{equation}
For fixed $q$, put $u_c(s)=q_ce^{sq_c}/\sum_dq_de^{sq_d}$.
Differentiation gives
\[
 \frac{d}{ds}\sum_cu_c(s)^2
 =\sum_{c,d}u_c(s)u_d(s)
       \big(u_c(s)-u_d(s)\big)(q_c-q_d)\ge0.
\]
The ordering of $u(s)$ is that of $q$, and the inequality is strict whenever
$q$ is nonuniform. Integrating from $0$ to $h_k$ proves the diversity claim.
Equation~\eqref{eq:finite-correct-tilt} also preserves the leader and makes
its conditional probability nondecreasing.

Every correct logit increases and every incorrect logit decreases, hence
$P_k$ increases. Jensen's inequality gives
\[
 \log\frac{P_{k+1}}{1-P_{k+1}}-\log\frac{P_k}{1-P_k}
 =\log\sum_cq_c e^{h_kq_c}-\log\sum_ar_a e^{-h_kr_a}
 \ge h_k\sum_cq_c^2\ge h_k/m.
\]
If $P_k$ had an interior limit, continuity and positivity of $c$ would imply
$\sum_kh_k=\infty$, contradicting bounded correctness odds. Thus $P_k\to1$.
Moreover, $\sum_kh_k=\infty$ is necessary: otherwise every logit would have
finite total variation and converge to a finite value, contradicting
$P_k\to1$.
Set $\alpha=q_{c^\star,0}$ and
$\xi_{c,k}=q_{c,k}/q_{c^\star,k}$. For $c\ne c^\star$,
\[
 \xi_{c,k+1}=\xi_{c,k}
  e^{-h_k(q_{c^\star,k}-q_{c,k})},
 \qquad
 q_{c^\star,k}-q_{c,k}
 \ge\alpha(1-\xi_{c,0})>0.
\]
The divergent sum of $h_k$ forces every such ratio to zero, proving
$q_{c^\star,k}\to1$.
\end{proof}

The update uses the exact expected gradient. Finite steps preserve the
concentration conclusion but can change the trajectory across estimators.
The identical diversity-versus-correctness curve in
Proposition~\ref{prop:pcmd-drift} requires infinitesimal mean flow;
it does not extend to these finite updates.

Finite multinomial counts can break a uniform policy's symmetry
\citep[Proposition~A.1]{lochab2026ucpo}. For the centered group estimator,
the leading expected one-step loss in \pmd{} is quadratic in the learning
rate, with a coefficient determined by group size and correctness.

\begin{proposition}[Sampled updates break conditional symmetry]
\enforcehalffinalproseline
\label{prop:stochastic-pcmd}
Retain (A1), (A2), (A4), and (A5) of
Assumption~\ref{ass:categorical-mean-flow}, and replace (A3) by one sampled
update $z^+=z+\eta\widehat g_G$ from
Equation~\eqref{eq:group-estimator}, with $\eta>0$. Suppose $0<P<1$ and $q_c=1/m$ for every correct mode. Then
\begin{equation}
 \E[D(q^+)]<1-1/m,
 \qquad \E[q_c^+]=1/m.
 \label{eq:stochastic-strict-decrease}
\end{equation}
As $\eta\to0$, for the fixed initial policy and estimator,
\begin{equation}
 \E[D(q^+)]
 =1-\frac1m
  -\eta^2\frac{m-1}{m^3}
    \E\!\left[\frac{w(R)^2R(G-R)^2}{G^4}\right]
  +O(\eta^3),
 \label{eq:stochastic-pcmd-expansion}
\end{equation}
where the summand is zero on $R=0,G$. For Dr.GRPO, $w=1$ and
\begin{equation}
 \E\!\left[\frac{R(G-R)^2}{G^4}\right]
 =\frac{G-1}{G^3}P(1-P)\big[1+(G-2)(1-P)\big].
 \label{eq:stochastic-dr-coefficient}
\end{equation}
\end{proposition}

\begin{proof}
\enforcehalffinalproseline
Let $N_c$ count occurrences of correct mode $c$. The centered advantages
sum to zero, so the common softmax-score term cancels and
$\widehat g_c=\kappa_RN_c$, where
$\kappa_R=w(R)(G-R)/G^2$ on mixed groups and $\kappa_R=0$ otherwise.
Conditional on $R$, the correct counts are
$N\sim\operatorname{Multinomial}(R;1/m,\ldots,1/m)$, and
\[
 q_c^+=\frac{e^{\eta\kappa_RN_c}}{\sum_de^{\eta\kappa_RN_d}}.
\]
Diversity is at most $1-1/m$, with equality only at uniform $q^+$.
The event $R=1$ has positive probability; on this event one correct logit
increases and the others do not, so $q^+$ is nonuniform. This proves the
strict inequality for every $\eta>0$. Permutation symmetry of the correct
labels gives $\E[q_c^+]=1/m$.

For a fixed vector $v$, Taylor expansion about uniform logits gives
\[
 \sum_c\operatorname{softmax}(\eta v)_c^2
 =\frac1m+\frac{\eta^2}{m^2}\sum_c(v_c-\bar v)^2+O(\eta^3),
 \qquad \bar v=\frac1m\sum_cv_c.
\]
Use $v_c=\kappa_RN_c$ and
$\E[\sum_c(N_c-R/m)^2\mid R]=R(1-1/m)$.
There are finitely many possible groups and all mixed-group weights are
finite, so the Taylor remainder is uniform over those groups. Averaging
proves Equation~\eqref{eq:stochastic-pcmd-expansion}.
Finally,
\[
 \E[R(G-R)^2]
 =G(G-1)(G-2)P(1-P)^2+G(G-1)P(1-P),
\]
by the binomial factorial moments, yielding
Equation~\eqref{eq:stochastic-dr-coefficient}.
\end{proof}

At fixed $P\in(0,1)$, the Dr.GRPO coefficient in
Equation~\eqref{eq:stochastic-dr-coefficient} is
$P(1-P)^2/G+O(G^{-2})$ as $G\to\infty$. The coefficient of the
quadratic learning-rate term therefore decays asymptotically as $1/G$;
it need not decrease at every increment in finite group size. This sampling
effect is absent from the uniform mean flow. The result concerns one step
from a uniform correct-mode distribution and gives neither monotone expected
diversity from an arbitrary initial policy nor stochastic-optimizer convergence.

\begin{corollary}[Finite-budget disappearance of competing modes]
\enforcehalffinalproseline
\label{cor:finite-budget-invisibility}
In the mean flow of Theorem~\ref{thm:grpo-collapse}, let $c^\star$ be the
unique initial correct leader and define
\[
 \alpha=q_{c^\star}(0),
 \qquad
 \Gamma_0=\sum_{c\ne c^\star}
             \frac{q_c(0)}{\alpha-q_c(0)}.
\]
At effective time $\tau$,
\begin{equation}
 1-q_{c^\star}(\tau)\le\Gamma_0e^{-\alpha\tau},
 \qquad D(q(\tau))\le2\Gamma_0e^{-\alpha\tau}.
 \label{eq:finite-minority-bound}
\end{equation}
For $K\ge1$ independent full-policy draws at that checkpoint,
\begin{equation}
 \Pr(\text{any correct mode other than }c^\star\text{ is observed})
 \le K\Gamma_0e^{-\alpha\tau}.
 \label{eq:finite-inference-visibility}
\end{equation}
In particular, if
$\tau\ge\max\{0,\alpha^{-1}\log(K\Gamma_0/\delta)\}$ for
$0<\delta<1$, no competing correct mode is observed with probability at
least $1-\delta$, although every mode retains positive probability.
\end{corollary}

\begin{proof}
\enforcehalffinalproseline
For $c\ne c^\star$, write $\xi_c=q_c/q_{c^\star}<1$.
The leader is nondecreasing, and Equation~\eqref{eq:grpo-replicator} gives
\[
 \frac{d}{d\tau}\log\frac{\xi_c}{1-\xi_c}
 =-q_{c^\star}\le-\alpha.
\]
Hence
$\xi_c(\tau)\le[q_c(0)/(\alpha-q_c(0))]e^{-\alpha\tau}$.
Now $1-q_{c^\star}=\sum_{c\ne c^\star}\xi_c/(1+\sum_{c\ne c^\star}\xi_c)$
and $D\le1-q_{c^\star}^2\le2(1-q_{c^\star})$, proving the first bounds.
A single full-policy draw hits a competing correct mode with probability
$P(1-q_{c^\star})\le1-q_{c^\star}$; a union bound over $K$ draws completes
the proof.
\end{proof}

The bound gives a finite-budget interpretation of the rare-output pruning
mechanism discussed by \citet[Section~4.2.5]{lochab2026ucpo}.
It can also be indexed by correctness: Proposition~\ref{prop:pcmd-drift}
and $S_2+T_2\le2$ imply
$\tau\ge\tfrac12\log\{[P/(1-P)]/[P(0)/(1-P(0))]\}$.
These are sufficient visibility bounds within the categorical model;
they do not require literal support loss or identify an empirical collapse
time for neural training.


\subsection{A replay signal when fresh groups supply none}
